In summary, our results demonstrate that 1) LLM alignment with the human brain is driven by both model size and performance on divergent thinking assessments, 2) alignment is stronger for later layers than earlier layers, and 3) post-training can increase this alignment, but only if the post-training objective is explicitly designed to optimize alignment with human behavior, or to push models towards greater diversity in responses. Notably, however, these effects are stage-dependent: the relationship between model properties and brain alignment is clearest during prompt processing, and weakens once models begin generating responses, suggesting that the representational dynamics of language generation diverge from human creative cognition in ways that are not simply resolved by scaling. Perhaps most strikingly, our post-training results reveal that the choice of training objective has interpretable and selective consequences for how model representations relate to the neural geometry of creative thought: reasoning-chain training steers representations away from high-creativity neural responses entirely, while creativity-optimized training produces the opposite effect. 

Our post-training results carry important implications for how language models should be fine-tuned in future work. Post-training strategies and datasets have thus far emphasized tasks requiring \textit{convergent thinking}, such as solving difficult math problems or coding tasks \citep{deepseekai2025deepseekr1incentivizingreasoningcapability, lightman2023let, chen2021evaluating}, because they are easily evaluated using a human-provided ground truth or an LLM judge. While this has proven effective for developing models that can solve tasks which emphasize convergent thinking --- reasoning toward a single correct answer or solution --- our findings indicate that it may simultaneously impair the ability of LLMs to come up with many and varied ideas, the hallmark of divergent thinking \citep{runco2012divergent}. Specifically, a model trained for reasoning consistently demonstrated stronger alignment for low-creativity responses than high-creativity ones, suggesting that its post-training had inadvertently shifted it toward neural representations associated with lower creativity. Together, these findings suggest that brain alignment offers a uniquely informative lens for evaluating and developing language models in the context of creative cognition — one that goes beyond behavioral benchmarks and captures something closer to the computational principles underlying human creative thought. Given the crucial role of creativity in fostering scientific, engineering, and artistic advances \citep{torrance1974torrance, guilford1967nature}, we must develop LLMs capable of creativity when tasked with solving challenging, long-horizon, or ambiguously scoped problems. Our findings suggest that post-training exclusively with convergent-thinking tasks may actively degrade models' capacity for divergent thinking.